\section*{M5b: Growing-tree convergence under weighted summability}

The fixed-tree theorem does not extend to $k_N\to\infty$ from pointwise
consistency alone. The following corrected statement isolates the missing
hypothesis in the form used by the multilinear telescoping proof.

Let $T_N$ be a sequence of finite typed projected trees. Let $\mathcal I_N$
be its internal vertices, let $\delta_{v,N}\ge0$ be the local
discretization/commutation defect at vertex $v$, and let $w_{v,N}\ge0$ be a
downstream gain bound for carrying a perturbation at $v$ to the observed
root. Assume the nodewise telescoping argument proves
\[
 \|J_NF_{T_N,N}-F_{T_N}\|+
 \|J_NR_{T_N,N}-R_{T_N}\|
 \le \sum_{v\in\mathcal I_N}w_{v,N}\delta_{v,N}.
 \tag{M5b.1}
\]
Assume also the weighted summability condition
\[
 \boxed{\quad\sum_{v\in\mathcal I_N}w_{v,N}\delta_{v,N}\longrightarrow0.\quad}
 \tag{M5b.2}
\]
No bound on $|\mathcal I_N|$ is required.

\subsection*{Theorem}

Under (M5b.1)--(M5b.2), the discretized ambient and recursively projected
evaluations converge at the observed root. If the projected-root bound holds
at every $N$ with right-hand side $B_N$, then
\[
 E_{T_N,N}^{\mathrm{proj}}\longrightarrow E_{T_N}^{\mathrm{proj}},
 \qquad
 E_{T_N}^{\mathrm{proj}}\le\limsup_{N\to\infty}B_N.
\]

\subsection*{Proof}

The first conclusion follows immediately from (M5b.1) and (M5b.2), since
both nonnegative root errors are bounded by a quantity tending to zero. The
projected-root error is a continuous norm of the difference between the
projected ambient and reduced root values, so the same convergence gives the
first limit. For every $N$, the finite-tree projected theorem gives
$E_{T_N,N}^{\mathrm{proj}}\le B_N$; passing to a subsequence along which
$B_N$ approaches its limsup preserves the non-strict inequality. The complete
growing sum is controlled before taking the limit, so the proof is independent
of the number of internal vertices. $\blacksquare$

\subsection*{Boundary}

This is a weighted summability condition, not pointwise convergence at each
fixed vertex. If $w_{v,N}\le W$ and $|\mathcal I_N|=k_N$, the sufficient rate
$\sup_v\delta_{v,N}=o(k_N^{-1})$ implies (M5b.2). The borderline rate
$\Theta(k_N^{-1})$ need not imply convergence, as explained by the existing
counterexample in
\texttt{counterexamples\_without\_uniformity.tex}.

\subsection*{Status}

\textsc{proved\_under\_stated\_assumptions}. This closes the corrected
summability formulation, not unrestricted growing-tree convergence.
